\documentclass[aspectratio=169,10pt]{beamer}
\makeatletter\def\input@path{{../../estilo/}{../estilo/}{estilo/}}\makeatother
\usepackage{devmobile}
\modulo[Flutter]{08}{Arquitetura e Gerência de Estado}{Arquitetura recomendada (MVVM), views, view models, repositórios, serviços, Result, Command, injeção de dependência, Provider, Riverpod e Bloc}

\begin{document}
\capa

\begin{frame}{Objetivos do módulo}
  \begin{conceito}[Ao final deste módulo você será capaz de]
    \begin{itemize}
      \item Aplicar a arquitetura recomendada pela documentação do Flutter: camada de UI e camada de dados.
      \item Implementar \emph{views}, \emph{view models}, repositórios e serviços.
      \item Tratar falhas com o tipo \cod{Result} e ações assíncronas com o padrão \textbf{Command}.
      \item Injetar dependências com \textbf{Provider}.
      \item Comparar Provider, Riverpod e Bloc e escolher a abordagem adequada.
      \item Organizar as pastas de um app que vai crescer.
    \end{itemize}
  \end{conceito}
\end{frame}

\sumario

% =============================================================================
%  Fundamentos compartilhados --- Módulo 08: arquitetura de software
%  Incluído pelas trilhas Kotlin e Flutter com % =============================================================================
%  Fundamentos compartilhados --- Módulo 08: arquitetura de software
%  Incluído pelas trilhas Kotlin e Flutter com % =============================================================================
%  Fundamentos compartilhados --- Módulo 08: arquitetura de software
%  Incluído pelas trilhas Kotlin e Flutter com \input{../fundamentos/f08-arquitetura}
%  (coesão, acoplamento, SOLID e padrões de projeto estão no Módulo 03)
% =============================================================================
\section{Fundamentos: arquitetura de software}

\begin{frame}{O que é arquitetura de software}
  \begin{columns}[T]
    \begin{column}{0.52\textwidth}
      \begin{definicao}[Arquitetura]
        As grandes estruturas de um sistema: seus elementos, as relações entre eles e as decisões \textbf{difíceis de mudar depois}. Ela existe para proteger atributos de qualidade.
      \end{definicao}
      \scriptsize
      \begin{tabularx}{\linewidth}{@{}l L@{}}
        \toprule
        \cabecalho{Atributo} & \cabecalho{Pergunta no LabMonitor} \\
        \midrule
        Manutenibilidade & trocar a API exige mexer nas telas? \\
        Testabilidade & dá para testar a regra de reserva sem um aparelho? \\
        Desempenho & a lista de leituras rola sem engasgar? \\
        Funcionamento \emph{offline} & o app mostra dados sem internet? \\
        Segurança & o \emph{token} fica protegido? \\
        \bottomrule
      \end{tabularx}
    \end{column}
    \begin{column}{0.46\textwidth}
      \begin{analogia}
        A planta de uma casa decide por onde passam os canos antes da primeira parede. Pintar um quarto não é arquitetura; mudar a cozinha de lugar é.
      \end{analogia}
      \begin{conceito}[Boa arquitetura em um app]
        Permite mudar a interface sem tocar nos dados, trocar a API sem tocar na interface e testar regras sem abrir o app.
      \end{conceito}
      \begin{dica}
        Coesão, acoplamento, SOLID e padrões de projeto foram vistos no Módulo 03: aqui eles ganham escala.
      \end{dica}
    \end{column}
  \end{columns}
\end{frame}

\begin{frame}{Padrões de apresentação: de MVC a MVI}
  \centering
  \begin{tikzpicture}[every node/.style={font=\tiny}]
    \draw[seta] (0,0) -- (13.4,0);
    \foreach \x/\t/\d in {1.2/MVC/{1979, Smalltalk\\ (Trygve Reenskaug)},4.6/MVP/{anos 1990 (Taligent);\\ comum no Android até 2017},8.1/MVVM/{2005, Microsoft (WPF);\\ padrão no Android e no Flutter},11.6/MVI/{anos 2010, inspirado\\ no Redux e no Cycle.js}}{
      \fill[ambar] (\x,0) circle (2pt);
      \node[font=\scriptsize\bfseries,text=azulpetroleo,above=1mm] at (\x,0) {\t};
      \node[rotulo,below=1mm] at (\x,0) {\d};
    }
  \end{tikzpicture}
  \raggedright
  \scriptsize
  \begin{tabularx}{\textwidth}{@{}l L L L@{}}
    \toprule
    \cabecalho{Padrão} & \cabecalho{Quem atualiza a tela} & \cabecalho{Ponto forte} & \cabecalho{Ponto fraco} \\
    \midrule
    MVC & o \emph{controller} manipula a \emph{view} & simples & no Android, a \cod{Activity} virava um \emph{controller} gigante \\
    MVP & o \emph{presenter} chama métodos de uma interface da \emph{view} & testável sem a tela & muito código de interface; guarda referência à \emph{view} \\
    MVVM & a \emph{view} observa o estado exposto pelo \emph{view model} & casa com UI declarativa e sobrevive à rotação & vários fluxos de estado podem se desencontrar \\
    MVI & estado único e imutável; cada intenção gera um novo estado & previsível e fácil de depurar & cerimônia demais para telas simples \\
    \bottomrule
  \end{tabularx}
  \normalsize
  \begin{conceito}[Na prática atual]
    Compose e Flutter usam MVVM com um estado único por tela (\cod{UiState}): uma mistura de MVVM com a disciplina do MVI.
  \end{conceito}
\end{frame}

\begin{frame}{Arquitetura limpa e arquitetura hexagonal}
  \begin{columns}[c]
    \begin{column}{0.46\textwidth}
      \centering
      \begin{tikzpicture}[every node/.style={font=\tiny,align=center}]
        \fill[cinzaclaro] (0,0) circle (2.15);
        \fill[cinzafundo,draw=azulpetroleo!40] (0,0) circle (1.6);
        \fill[ambarclaro!60,draw=ambar] (0,0) circle (1.05);
        \fill[azulpetroleo] (0,0) circle (0.55);
        \node[text=white,font=\tiny\bfseries] at (0,0) {Entidades};
        \node at (0,0.8) {Casos de uso};
        \node at (0,1.33) {Adaptadores};
        \node at (0,1.88) {Frameworks e drivers};
        \node[rotulo] at (0,-1.33) {\emph{view models}, repositórios};
        \node[rotulo] at (0,-1.88) {UI, banco, HTTP, BLE};
        \draw[setaambar] (2.6,-1.2) -- (1.2,-0.55);
        \node[rotulo,text width=1.5cm] at (3.05,-1.5) {dependências apontam para dentro};
      \end{tikzpicture}
    \end{column}
    \begin{column}{0.52\textwidth}
      \begin{definicao}[Regra da dependência]
        O código das camadas internas não conhece as externas: as regras de negócio não importam nada de Android, Flutter, HTTP ou SQL.
      \end{definicao}
      \begin{historia}
        Alistair Cockburn propôs a arquitetura hexagonal, de portas e adaptadores, em 2005. Robert C. Martin reuniu ideias parecidas na Arquitetura Limpa, em 2012.
      \end{historia}
      \begin{dica}[Na medida certa]
        Em apps pequenos, camadas de UI e de dados bastam. Crie a camada de domínio quando regras se repetirem ou precisarem de testes isolados, como recomendam o Android e o Flutter.
      \end{dica}
    \end{column}
  \end{columns}
\end{frame}

\begin{frame}[fragile]{Inversão de dependência na prática}
  \begin{columns}[T]
    \begin{column}{0.5\textwidth}
      \centering
      \begin{tikzpicture}[every node/.style={font=\tiny}]
        \node[rotulo,anchor=west] at (-1.1,0.55) {\textbf{sem inversão}};
        \node[caixaescura,font=\tiny,minimum width=2cm] (r1) at (0,0) {regras de leitura};
        \node[caixa,font=\tiny,minimum width=1.8cm] (b1) at (3.4,0) {\cod{FonteBle}};
        \draw[seta] (r1) -- node[rotulo,above]{conhece} (b1);
        \node[rotulo,anchor=west] at (-1.1,-0.75) {\textbf{com inversão}};
        \node[caixaescura,font=\tiny,minimum width=2cm] (r2) at (0,-1.3) {regras de leitura};
        \node[caixaambar,font=\tiny,minimum width=2.4cm] (i2) at (3.4,-1.3) {\emph{interface}\\ \cod{FonteLeituras}};
        \node[caixa,font=\tiny,minimum width=1.35cm] (b2) at (1.95,-2.75) {\cod{FonteBle}};
        \node[caixa,font=\tiny,minimum width=1.35cm] (m2) at (3.4,-2.75) {\cod{FonteMqtt}};
        \node[caixa,font=\tiny,minimum width=1.35cm] (f2) at (4.85,-2.75) {\cod{FonteFake}};
        \draw[seta] (r2) -- node[rotulo,above]{usa} (i2);
        \foreach \n in {b2,m2,f2}{\draw[setaambar] (\n.north) -- (\n.north |- i2.south);}
        \node[rotulo,below=1mm of m2] {implementam a interface};
      \end{tikzpicture}
\begin{codigo}[title=Dart]{Dart}
abstract interface class FonteLeituras {
  Stream<Leitura> leituras();
}
\end{codigo}
    \end{column}
    \begin{column}{0.48\textwidth}
      \begin{definicao}[Inversão de dependência]
        Módulos de alto nível não dependem de detalhes; ambos dependem de abstrações. Quem precisa define a interface; a camada de dados a implementa.
      \end{definicao}
\begin{codigo}[title=Kotlin]{Kotlin}
interface FonteLeituras {
  fun leituras(): Flow<Leitura>
}
\end{codigo}
      \begin{conceito}[O que se ganha]
        Trocar BLE por MQTT sem mexer nas regras, testar com uma fonte falsa e fazer a tela antes de o hardware existir.
      \end{conceito}
    \end{column}
  \end{columns}
\end{frame}

\begin{frame}{Injeção de dependência e localizador de serviços}
  \begin{columns}[T]
    \begin{column}{0.5\textwidth}
      \begin{definicao}[Inversão de controle]
        A classe não cria o que usa: alguém de fora monta os objetos e os entrega. É o princípio de Hollywood: ``não nos chame, nós chamamos você''.
      \end{definicao}
      \centering
      \begin{tikzpicture}[every node/.style={font=\tiny}]
        \node[caixaambarforte,font=\tiny\bfseries,minimum width=5.6cm] (raiz) at (0,0) {raiz de composição (\cod{Application}, \cod{main})};
        \node[caixa,font=\tiny,minimum width=1.7cm] (h) at (-1.9,-1) {cliente HTTP};
        \node[caixa,font=\tiny,minimum width=1.7cm] (b) at (0,-1) {banco};
        \node[caixapetroleo,font=\tiny,minimum width=1.9cm] (r) at (-0.95,-2) {repositório};
        \node[caixaescura,font=\tiny,minimum width=1.7cm] (v) at (1.9,-2) {\emph{view model}};
        \draw[setatrac] (raiz.south -| h) -- (h); \draw[setatrac] (raiz.south -| b) -- (b);
        \draw[seta] (h) -- (r); \draw[seta] (b) -- (r);
        \draw[seta] (r) -- node[rotulo,above]{injeta} (v);
        \node[rotulo] at (1.9,-1.05) {tracejado: cria};
      \end{tikzpicture}
      \raggedright
      \begin{historia}
        Martin Fowler batizou o padrão de injeção de dependência em 2004.
      \end{historia}
    \end{column}
    \begin{column}{0.48\textwidth}
      \scriptsize
      \begin{tabularx}{\linewidth}{@{}l L@{}}
        \toprule
        \cabecalho{Abordagem} & \cabecalho{Como funciona} \\
        \midrule
        Construtor & recebe tudo pronto; dependências explícitas e fáceis de trocar em testes \\
        Contêiner & uma biblioteca monta o grafo: Hilt, Koin; Provider, Riverpod \\
        Localizador & a classe pede a um registro global (\cod{get\_it}): prático, mas esconde dependências e só falha em execução \\
        \bottomrule
      \end{tabularx}
      \normalsize
      \begin{pordentro}[Tempo de vida (escopo)]
        \textbf{Único no app}: cliente HTTP, banco. \textbf{Por tela}: \emph{view model}. \textbf{Por uso}: objeto novo a cada pedido. Escopo errado causa vazamento ou estado compartilhado por engano.
      \end{pordentro}
    \end{column}
  \end{columns}
\end{frame}

\begin{frame}{Decisões de arquitetura: trocas e registros}
  \begin{columns}[T]
    \begin{column}{0.5\textwidth}
      \begin{conceito}[Toda decisão é uma troca]
        Mais camadas trazem isolamento e testabilidade, mas também mais arquivos e indireção. Pese o tamanho do time, a vida esperada do app e o custo de errar.
      \end{conceito}
      \scriptsize
      \begin{tabularx}{\linewidth}{@{}l L@{}}
        \toprule
        \cabecalho{Princípio} & \cabecalho{Significado} \\
        \midrule
        YAGNI & não construa o que ainda não é necessário \\
        KISS & prefira a solução mais simples que funciona \\
        DRY & cada conhecimento em um só lugar \\
        Regra de três & abstraia na terceira repetição, não na primeira \\
        \bottomrule
      \end{tabularx}
    \end{column}
    \begin{column}{0.48\textwidth}
      \begin{definicao}[ADR (\emph{Architecture Decision Record})]
        Texto curto, versionado com o código, que registra uma decisão: contexto, opções, decisão e consequências. Responde no futuro: ``por que fizemos assim?''.
      \end{definicao}
      \begin{exemplo}[ADR 007 --- Estado das telas]
        \textbf{Contexto}: leituras do sensor em tempo real. \textbf{Decisão}: um \cod{UiState} imutável por tela, exposto como fluxo observável. \textbf{Consequência}: telas testáveis sem interface.
      \end{exemplo}
      \begin{historia}
        Michael Nygard propôs os ADRs em 2011.
      \end{historia}
    \end{column}
  \end{columns}
\end{frame}

%  (coesão, acoplamento, SOLID e padrões de projeto estão no Módulo 03)
% =============================================================================
\section{Fundamentos: arquitetura de software}

\begin{frame}{O que é arquitetura de software}
  \begin{columns}[T]
    \begin{column}{0.52\textwidth}
      \begin{definicao}[Arquitetura]
        As grandes estruturas de um sistema: seus elementos, as relações entre eles e as decisões \textbf{difíceis de mudar depois}. Ela existe para proteger atributos de qualidade.
      \end{definicao}
      \scriptsize
      \begin{tabularx}{\linewidth}{@{}l L@{}}
        \toprule
        \cabecalho{Atributo} & \cabecalho{Pergunta no LabMonitor} \\
        \midrule
        Manutenibilidade & trocar a API exige mexer nas telas? \\
        Testabilidade & dá para testar a regra de reserva sem um aparelho? \\
        Desempenho & a lista de leituras rola sem engasgar? \\
        Funcionamento \emph{offline} & o app mostra dados sem internet? \\
        Segurança & o \emph{token} fica protegido? \\
        \bottomrule
      \end{tabularx}
    \end{column}
    \begin{column}{0.46\textwidth}
      \begin{analogia}
        A planta de uma casa decide por onde passam os canos antes da primeira parede. Pintar um quarto não é arquitetura; mudar a cozinha de lugar é.
      \end{analogia}
      \begin{conceito}[Boa arquitetura em um app]
        Permite mudar a interface sem tocar nos dados, trocar a API sem tocar na interface e testar regras sem abrir o app.
      \end{conceito}
      \begin{dica}
        Coesão, acoplamento, SOLID e padrões de projeto foram vistos no Módulo 03: aqui eles ganham escala.
      \end{dica}
    \end{column}
  \end{columns}
\end{frame}

\begin{frame}{Padrões de apresentação: de MVC a MVI}
  \centering
  \begin{tikzpicture}[every node/.style={font=\tiny}]
    \draw[seta] (0,0) -- (13.4,0);
    \foreach \x/\t/\d in {1.2/MVC/{1979, Smalltalk\\ (Trygve Reenskaug)},4.6/MVP/{anos 1990 (Taligent);\\ comum no Android até 2017},8.1/MVVM/{2005, Microsoft (WPF);\\ padrão no Android e no Flutter},11.6/MVI/{anos 2010, inspirado\\ no Redux e no Cycle.js}}{
      \fill[ambar] (\x,0) circle (2pt);
      \node[font=\scriptsize\bfseries,text=azulpetroleo,above=1mm] at (\x,0) {\t};
      \node[rotulo,below=1mm] at (\x,0) {\d};
    }
  \end{tikzpicture}
  \raggedright
  \scriptsize
  \begin{tabularx}{\textwidth}{@{}l L L L@{}}
    \toprule
    \cabecalho{Padrão} & \cabecalho{Quem atualiza a tela} & \cabecalho{Ponto forte} & \cabecalho{Ponto fraco} \\
    \midrule
    MVC & o \emph{controller} manipula a \emph{view} & simples & no Android, a \cod{Activity} virava um \emph{controller} gigante \\
    MVP & o \emph{presenter} chama métodos de uma interface da \emph{view} & testável sem a tela & muito código de interface; guarda referência à \emph{view} \\
    MVVM & a \emph{view} observa o estado exposto pelo \emph{view model} & casa com UI declarativa e sobrevive à rotação & vários fluxos de estado podem se desencontrar \\
    MVI & estado único e imutável; cada intenção gera um novo estado & previsível e fácil de depurar & cerimônia demais para telas simples \\
    \bottomrule
  \end{tabularx}
  \normalsize
  \begin{conceito}[Na prática atual]
    Compose e Flutter usam MVVM com um estado único por tela (\cod{UiState}): uma mistura de MVVM com a disciplina do MVI.
  \end{conceito}
\end{frame}

\begin{frame}{Arquitetura limpa e arquitetura hexagonal}
  \begin{columns}[c]
    \begin{column}{0.46\textwidth}
      \centering
      \begin{tikzpicture}[every node/.style={font=\tiny,align=center}]
        \fill[cinzaclaro] (0,0) circle (2.15);
        \fill[cinzafundo,draw=azulpetroleo!40] (0,0) circle (1.6);
        \fill[ambarclaro!60,draw=ambar] (0,0) circle (1.05);
        \fill[azulpetroleo] (0,0) circle (0.55);
        \node[text=white,font=\tiny\bfseries] at (0,0) {Entidades};
        \node at (0,0.8) {Casos de uso};
        \node at (0,1.33) {Adaptadores};
        \node at (0,1.88) {Frameworks e drivers};
        \node[rotulo] at (0,-1.33) {\emph{view models}, repositórios};
        \node[rotulo] at (0,-1.88) {UI, banco, HTTP, BLE};
        \draw[setaambar] (2.6,-1.2) -- (1.2,-0.55);
        \node[rotulo,text width=1.5cm] at (3.05,-1.5) {dependências apontam para dentro};
      \end{tikzpicture}
    \end{column}
    \begin{column}{0.52\textwidth}
      \begin{definicao}[Regra da dependência]
        O código das camadas internas não conhece as externas: as regras de negócio não importam nada de Android, Flutter, HTTP ou SQL.
      \end{definicao}
      \begin{historia}
        Alistair Cockburn propôs a arquitetura hexagonal, de portas e adaptadores, em 2005. Robert C. Martin reuniu ideias parecidas na Arquitetura Limpa, em 2012.
      \end{historia}
      \begin{dica}[Na medida certa]
        Em apps pequenos, camadas de UI e de dados bastam. Crie a camada de domínio quando regras se repetirem ou precisarem de testes isolados, como recomendam o Android e o Flutter.
      \end{dica}
    \end{column}
  \end{columns}
\end{frame}

\begin{frame}[fragile]{Inversão de dependência na prática}
  \begin{columns}[T]
    \begin{column}{0.5\textwidth}
      \centering
      \begin{tikzpicture}[every node/.style={font=\tiny}]
        \node[rotulo,anchor=west] at (-1.1,0.55) {\textbf{sem inversão}};
        \node[caixaescura,font=\tiny,minimum width=2cm] (r1) at (0,0) {regras de leitura};
        \node[caixa,font=\tiny,minimum width=1.8cm] (b1) at (3.4,0) {\cod{FonteBle}};
        \draw[seta] (r1) -- node[rotulo,above]{conhece} (b1);
        \node[rotulo,anchor=west] at (-1.1,-0.75) {\textbf{com inversão}};
        \node[caixaescura,font=\tiny,minimum width=2cm] (r2) at (0,-1.3) {regras de leitura};
        \node[caixaambar,font=\tiny,minimum width=2.4cm] (i2) at (3.4,-1.3) {\emph{interface}\\ \cod{FonteLeituras}};
        \node[caixa,font=\tiny,minimum width=1.35cm] (b2) at (1.95,-2.75) {\cod{FonteBle}};
        \node[caixa,font=\tiny,minimum width=1.35cm] (m2) at (3.4,-2.75) {\cod{FonteMqtt}};
        \node[caixa,font=\tiny,minimum width=1.35cm] (f2) at (4.85,-2.75) {\cod{FonteFake}};
        \draw[seta] (r2) -- node[rotulo,above]{usa} (i2);
        \foreach \n in {b2,m2,f2}{\draw[setaambar] (\n.north) -- (\n.north |- i2.south);}
        \node[rotulo,below=1mm of m2] {implementam a interface};
      \end{tikzpicture}
\begin{codigo}[title=Dart]{Dart}
abstract interface class FonteLeituras {
  Stream<Leitura> leituras();
}
\end{codigo}
    \end{column}
    \begin{column}{0.48\textwidth}
      \begin{definicao}[Inversão de dependência]
        Módulos de alto nível não dependem de detalhes; ambos dependem de abstrações. Quem precisa define a interface; a camada de dados a implementa.
      \end{definicao}
\begin{codigo}[title=Kotlin]{Kotlin}
interface FonteLeituras {
  fun leituras(): Flow<Leitura>
}
\end{codigo}
      \begin{conceito}[O que se ganha]
        Trocar BLE por MQTT sem mexer nas regras, testar com uma fonte falsa e fazer a tela antes de o hardware existir.
      \end{conceito}
    \end{column}
  \end{columns}
\end{frame}

\begin{frame}{Injeção de dependência e localizador de serviços}
  \begin{columns}[T]
    \begin{column}{0.5\textwidth}
      \begin{definicao}[Inversão de controle]
        A classe não cria o que usa: alguém de fora monta os objetos e os entrega. É o princípio de Hollywood: ``não nos chame, nós chamamos você''.
      \end{definicao}
      \centering
      \begin{tikzpicture}[every node/.style={font=\tiny}]
        \node[caixaambarforte,font=\tiny\bfseries,minimum width=5.6cm] (raiz) at (0,0) {raiz de composição (\cod{Application}, \cod{main})};
        \node[caixa,font=\tiny,minimum width=1.7cm] (h) at (-1.9,-1) {cliente HTTP};
        \node[caixa,font=\tiny,minimum width=1.7cm] (b) at (0,-1) {banco};
        \node[caixapetroleo,font=\tiny,minimum width=1.9cm] (r) at (-0.95,-2) {repositório};
        \node[caixaescura,font=\tiny,minimum width=1.7cm] (v) at (1.9,-2) {\emph{view model}};
        \draw[setatrac] (raiz.south -| h) -- (h); \draw[setatrac] (raiz.south -| b) -- (b);
        \draw[seta] (h) -- (r); \draw[seta] (b) -- (r);
        \draw[seta] (r) -- node[rotulo,above]{injeta} (v);
        \node[rotulo] at (1.9,-1.05) {tracejado: cria};
      \end{tikzpicture}
      \raggedright
      \begin{historia}
        Martin Fowler batizou o padrão de injeção de dependência em 2004.
      \end{historia}
    \end{column}
    \begin{column}{0.48\textwidth}
      \scriptsize
      \begin{tabularx}{\linewidth}{@{}l L@{}}
        \toprule
        \cabecalho{Abordagem} & \cabecalho{Como funciona} \\
        \midrule
        Construtor & recebe tudo pronto; dependências explícitas e fáceis de trocar em testes \\
        Contêiner & uma biblioteca monta o grafo: Hilt, Koin; Provider, Riverpod \\
        Localizador & a classe pede a um registro global (\cod{get\_it}): prático, mas esconde dependências e só falha em execução \\
        \bottomrule
      \end{tabularx}
      \normalsize
      \begin{pordentro}[Tempo de vida (escopo)]
        \textbf{Único no app}: cliente HTTP, banco. \textbf{Por tela}: \emph{view model}. \textbf{Por uso}: objeto novo a cada pedido. Escopo errado causa vazamento ou estado compartilhado por engano.
      \end{pordentro}
    \end{column}
  \end{columns}
\end{frame}

\begin{frame}{Decisões de arquitetura: trocas e registros}
  \begin{columns}[T]
    \begin{column}{0.5\textwidth}
      \begin{conceito}[Toda decisão é uma troca]
        Mais camadas trazem isolamento e testabilidade, mas também mais arquivos e indireção. Pese o tamanho do time, a vida esperada do app e o custo de errar.
      \end{conceito}
      \scriptsize
      \begin{tabularx}{\linewidth}{@{}l L@{}}
        \toprule
        \cabecalho{Princípio} & \cabecalho{Significado} \\
        \midrule
        YAGNI & não construa o que ainda não é necessário \\
        KISS & prefira a solução mais simples que funciona \\
        DRY & cada conhecimento em um só lugar \\
        Regra de três & abstraia na terceira repetição, não na primeira \\
        \bottomrule
      \end{tabularx}
    \end{column}
    \begin{column}{0.48\textwidth}
      \begin{definicao}[ADR (\emph{Architecture Decision Record})]
        Texto curto, versionado com o código, que registra uma decisão: contexto, opções, decisão e consequências. Responde no futuro: ``por que fizemos assim?''.
      \end{definicao}
      \begin{exemplo}[ADR 007 --- Estado das telas]
        \textbf{Contexto}: leituras do sensor em tempo real. \textbf{Decisão}: um \cod{UiState} imutável por tela, exposto como fluxo observável. \textbf{Consequência}: telas testáveis sem interface.
      \end{exemplo}
      \begin{historia}
        Michael Nygard propôs os ADRs em 2011.
      \end{historia}
    \end{column}
  \end{columns}
\end{frame}

%  (coesão, acoplamento, SOLID e padrões de projeto estão no Módulo 03)
% =============================================================================
\section{Fundamentos: arquitetura de software}

\begin{frame}{O que é arquitetura de software}
  \begin{columns}[T]
    \begin{column}{0.52\textwidth}
      \begin{definicao}[Arquitetura]
        As grandes estruturas de um sistema: seus elementos, as relações entre eles e as decisões \textbf{difíceis de mudar depois}. Ela existe para proteger atributos de qualidade.
      \end{definicao}
      \scriptsize
      \begin{tabularx}{\linewidth}{@{}l L@{}}
        \toprule
        \cabecalho{Atributo} & \cabecalho{Pergunta no LabMonitor} \\
        \midrule
        Manutenibilidade & trocar a API exige mexer nas telas? \\
        Testabilidade & dá para testar a regra de reserva sem um aparelho? \\
        Desempenho & a lista de leituras rola sem engasgar? \\
        Funcionamento \emph{offline} & o app mostra dados sem internet? \\
        Segurança & o \emph{token} fica protegido? \\
        \bottomrule
      \end{tabularx}
    \end{column}
    \begin{column}{0.46\textwidth}
      \begin{analogia}
        A planta de uma casa decide por onde passam os canos antes da primeira parede. Pintar um quarto não é arquitetura; mudar a cozinha de lugar é.
      \end{analogia}
      \begin{conceito}[Boa arquitetura em um app]
        Permite mudar a interface sem tocar nos dados, trocar a API sem tocar na interface e testar regras sem abrir o app.
      \end{conceito}
      \begin{dica}
        Coesão, acoplamento, SOLID e padrões de projeto foram vistos no Módulo 03: aqui eles ganham escala.
      \end{dica}
    \end{column}
  \end{columns}
\end{frame}

\begin{frame}{Padrões de apresentação: de MVC a MVI}
  \centering
  \begin{tikzpicture}[every node/.style={font=\tiny}]
    \draw[seta] (0,0) -- (13.4,0);
    \foreach \x/\t/\d in {1.2/MVC/{1979, Smalltalk\\ (Trygve Reenskaug)},4.6/MVP/{anos 1990 (Taligent);\\ comum no Android até 2017},8.1/MVVM/{2005, Microsoft (WPF);\\ padrão no Android e no Flutter},11.6/MVI/{anos 2010, inspirado\\ no Redux e no Cycle.js}}{
      \fill[ambar] (\x,0) circle (2pt);
      \node[font=\scriptsize\bfseries,text=azulpetroleo,above=1mm] at (\x,0) {\t};
      \node[rotulo,below=1mm] at (\x,0) {\d};
    }
  \end{tikzpicture}
  \raggedright
  \scriptsize
  \begin{tabularx}{\textwidth}{@{}l L L L@{}}
    \toprule
    \cabecalho{Padrão} & \cabecalho{Quem atualiza a tela} & \cabecalho{Ponto forte} & \cabecalho{Ponto fraco} \\
    \midrule
    MVC & o \emph{controller} manipula a \emph{view} & simples & no Android, a \cod{Activity} virava um \emph{controller} gigante \\
    MVP & o \emph{presenter} chama métodos de uma interface da \emph{view} & testável sem a tela & muito código de interface; guarda referência à \emph{view} \\
    MVVM & a \emph{view} observa o estado exposto pelo \emph{view model} & casa com UI declarativa e sobrevive à rotação & vários fluxos de estado podem se desencontrar \\
    MVI & estado único e imutável; cada intenção gera um novo estado & previsível e fácil de depurar & cerimônia demais para telas simples \\
    \bottomrule
  \end{tabularx}
  \normalsize
  \begin{conceito}[Na prática atual]
    Compose e Flutter usam MVVM com um estado único por tela (\cod{UiState}): uma mistura de MVVM com a disciplina do MVI.
  \end{conceito}
\end{frame}

\begin{frame}{Arquitetura limpa e arquitetura hexagonal}
  \begin{columns}[c]
    \begin{column}{0.46\textwidth}
      \centering
      \begin{tikzpicture}[every node/.style={font=\tiny,align=center}]
        \fill[cinzaclaro] (0,0) circle (2.15);
        \fill[cinzafundo,draw=azulpetroleo!40] (0,0) circle (1.6);
        \fill[ambarclaro!60,draw=ambar] (0,0) circle (1.05);
        \fill[azulpetroleo] (0,0) circle (0.55);
        \node[text=white,font=\tiny\bfseries] at (0,0) {Entidades};
        \node at (0,0.8) {Casos de uso};
        \node at (0,1.33) {Adaptadores};
        \node at (0,1.88) {Frameworks e drivers};
        \node[rotulo] at (0,-1.33) {\emph{view models}, repositórios};
        \node[rotulo] at (0,-1.88) {UI, banco, HTTP, BLE};
        \draw[setaambar] (2.6,-1.2) -- (1.2,-0.55);
        \node[rotulo,text width=1.5cm] at (3.05,-1.5) {dependências apontam para dentro};
      \end{tikzpicture}
    \end{column}
    \begin{column}{0.52\textwidth}
      \begin{definicao}[Regra da dependência]
        O código das camadas internas não conhece as externas: as regras de negócio não importam nada de Android, Flutter, HTTP ou SQL.
      \end{definicao}
      \begin{historia}
        Alistair Cockburn propôs a arquitetura hexagonal, de portas e adaptadores, em 2005. Robert C. Martin reuniu ideias parecidas na Arquitetura Limpa, em 2012.
      \end{historia}
      \begin{dica}[Na medida certa]
        Em apps pequenos, camadas de UI e de dados bastam. Crie a camada de domínio quando regras se repetirem ou precisarem de testes isolados, como recomendam o Android e o Flutter.
      \end{dica}
    \end{column}
  \end{columns}
\end{frame}

\begin{frame}[fragile]{Inversão de dependência na prática}
  \begin{columns}[T]
    \begin{column}{0.5\textwidth}
      \centering
      \begin{tikzpicture}[every node/.style={font=\tiny}]
        \node[rotulo,anchor=west] at (-1.1,0.55) {\textbf{sem inversão}};
        \node[caixaescura,font=\tiny,minimum width=2cm] (r1) at (0,0) {regras de leitura};
        \node[caixa,font=\tiny,minimum width=1.8cm] (b1) at (3.4,0) {\cod{FonteBle}};
        \draw[seta] (r1) -- node[rotulo,above]{conhece} (b1);
        \node[rotulo,anchor=west] at (-1.1,-0.75) {\textbf{com inversão}};
        \node[caixaescura,font=\tiny,minimum width=2cm] (r2) at (0,-1.3) {regras de leitura};
        \node[caixaambar,font=\tiny,minimum width=2.4cm] (i2) at (3.4,-1.3) {\emph{interface}\\ \cod{FonteLeituras}};
        \node[caixa,font=\tiny,minimum width=1.35cm] (b2) at (1.95,-2.75) {\cod{FonteBle}};
        \node[caixa,font=\tiny,minimum width=1.35cm] (m2) at (3.4,-2.75) {\cod{FonteMqtt}};
        \node[caixa,font=\tiny,minimum width=1.35cm] (f2) at (4.85,-2.75) {\cod{FonteFake}};
        \draw[seta] (r2) -- node[rotulo,above]{usa} (i2);
        \foreach \n in {b2,m2,f2}{\draw[setaambar] (\n.north) -- (\n.north |- i2.south);}
        \node[rotulo,below=1mm of m2] {implementam a interface};
      \end{tikzpicture}
\begin{codigo}[title=Dart]{Dart}
abstract interface class FonteLeituras {
  Stream<Leitura> leituras();
}
\end{codigo}
    \end{column}
    \begin{column}{0.48\textwidth}
      \begin{definicao}[Inversão de dependência]
        Módulos de alto nível não dependem de detalhes; ambos dependem de abstrações. Quem precisa define a interface; a camada de dados a implementa.
      \end{definicao}
\begin{codigo}[title=Kotlin]{Kotlin}
interface FonteLeituras {
  fun leituras(): Flow<Leitura>
}
\end{codigo}
      \begin{conceito}[O que se ganha]
        Trocar BLE por MQTT sem mexer nas regras, testar com uma fonte falsa e fazer a tela antes de o hardware existir.
      \end{conceito}
    \end{column}
  \end{columns}
\end{frame}

\begin{frame}{Injeção de dependência e localizador de serviços}
  \begin{columns}[T]
    \begin{column}{0.5\textwidth}
      \begin{definicao}[Inversão de controle]
        A classe não cria o que usa: alguém de fora monta os objetos e os entrega. É o princípio de Hollywood: ``não nos chame, nós chamamos você''.
      \end{definicao}
      \centering
      \begin{tikzpicture}[every node/.style={font=\tiny}]
        \node[caixaambarforte,font=\tiny\bfseries,minimum width=5.6cm] (raiz) at (0,0) {raiz de composição (\cod{Application}, \cod{main})};
        \node[caixa,font=\tiny,minimum width=1.7cm] (h) at (-1.9,-1) {cliente HTTP};
        \node[caixa,font=\tiny,minimum width=1.7cm] (b) at (0,-1) {banco};
        \node[caixapetroleo,font=\tiny,minimum width=1.9cm] (r) at (-0.95,-2) {repositório};
        \node[caixaescura,font=\tiny,minimum width=1.7cm] (v) at (1.9,-2) {\emph{view model}};
        \draw[setatrac] (raiz.south -| h) -- (h); \draw[setatrac] (raiz.south -| b) -- (b);
        \draw[seta] (h) -- (r); \draw[seta] (b) -- (r);
        \draw[seta] (r) -- node[rotulo,above]{injeta} (v);
        \node[rotulo] at (1.9,-1.05) {tracejado: cria};
      \end{tikzpicture}
      \raggedright
      \begin{historia}
        Martin Fowler batizou o padrão de injeção de dependência em 2004.
      \end{historia}
    \end{column}
    \begin{column}{0.48\textwidth}
      \scriptsize
      \begin{tabularx}{\linewidth}{@{}l L@{}}
        \toprule
        \cabecalho{Abordagem} & \cabecalho{Como funciona} \\
        \midrule
        Construtor & recebe tudo pronto; dependências explícitas e fáceis de trocar em testes \\
        Contêiner & uma biblioteca monta o grafo: Hilt, Koin; Provider, Riverpod \\
        Localizador & a classe pede a um registro global (\cod{get\_it}): prático, mas esconde dependências e só falha em execução \\
        \bottomrule
      \end{tabularx}
      \normalsize
      \begin{pordentro}[Tempo de vida (escopo)]
        \textbf{Único no app}: cliente HTTP, banco. \textbf{Por tela}: \emph{view model}. \textbf{Por uso}: objeto novo a cada pedido. Escopo errado causa vazamento ou estado compartilhado por engano.
      \end{pordentro}
    \end{column}
  \end{columns}
\end{frame}

\begin{frame}{Decisões de arquitetura: trocas e registros}
  \begin{columns}[T]
    \begin{column}{0.5\textwidth}
      \begin{conceito}[Toda decisão é uma troca]
        Mais camadas trazem isolamento e testabilidade, mas também mais arquivos e indireção. Pese o tamanho do time, a vida esperada do app e o custo de errar.
      \end{conceito}
      \scriptsize
      \begin{tabularx}{\linewidth}{@{}l L@{}}
        \toprule
        \cabecalho{Princípio} & \cabecalho{Significado} \\
        \midrule
        YAGNI & não construa o que ainda não é necessário \\
        KISS & prefira a solução mais simples que funciona \\
        DRY & cada conhecimento em um só lugar \\
        Regra de três & abstraia na terceira repetição, não na primeira \\
        \bottomrule
      \end{tabularx}
    \end{column}
    \begin{column}{0.48\textwidth}
      \begin{definicao}[ADR (\emph{Architecture Decision Record})]
        Texto curto, versionado com o código, que registra uma decisão: contexto, opções, decisão e consequências. Responde no futuro: ``por que fizemos assim?''.
      \end{definicao}
      \begin{exemplo}[ADR 007 --- Estado das telas]
        \textbf{Contexto}: leituras do sensor em tempo real. \textbf{Decisão}: um \cod{UiState} imutável por tela, exposto como fluxo observável. \textbf{Consequência}: telas testáveis sem interface.
      \end{exemplo}
      \begin{historia}
        Michael Nygard propôs os ADRs em 2011.
      \end{historia}
    \end{column}
  \end{columns}
\end{frame}


% =============================================================================
\section{A arquitetura recomendada}

\begin{frame}{Por que arquitetura?}
  \begin{columns}[c]
    \begin{column}{0.46\textwidth}
      \centering
      \begin{tikzpicture}
        \node[caixaescura,minimum width=4.4cm,minimum height=3.4cm] (a) {};
        \node[font=\small\bfseries,text=white] at ([yshift=-3mm]a.north) {TelaPainel (1\,500 linhas)};
        \foreach \t/\y in {widgets e layout/0.75,chamadas HTTP/0.3,SQL/-0.15,regras de negócio/-0.6,Bluetooth/-1.05}{
          \node[caixaambar,minimum width=3.4cm,minimum height=4mm,font=\tiny] at (0,\y) {\t};}
      \end{tikzpicture}
    \end{column}
    \begin{column}{0.52\textwidth}
      \begin{erro}[Tudo dentro do widget]
        \begin{itemize}
          \item Impossível testar regras sem montar a tela.
          \item Uma mudança na API quebra várias telas.
          \item Dados duplicados e inconsistentes entre telas.
        \end{itemize}
      \end{erro}
      \begin{conceito}[Princípios]
        \begin{itemize}
          \item Separação de responsabilidades.
          \item Fonte única da verdade.
          \item Fluxo unidirecional de dados.
          \item UI como função do estado.
        \end{itemize}
      \end{conceito}
    \end{column}
  \end{columns}
\end{frame}

\begin{frame}{Camadas: UI e dados}
  \begin{columns}[c]
    \begin{column}{0.58\textwidth}
      \centering
      \begin{tikzpicture}[node distance=3mm]
        \node[caixa,minimum width=2.6cm,minimum height=8mm] (v) {\textbf{View}\\ \tiny widgets};
        \node[caixaescura,minimum width=2.6cm,minimum height=8mm,right=6mm of v] (vm) {\textbf{View model}\\ \tiny \cod{ChangeNotifier}};
        \node[pontilhado,fit=(v)(vm),label={[rotulo,anchor=south west]north west:\textbf{Camada de UI}}] (ui) {};
        \node[caixapetroleo,minimum width=5.8cm,minimum height=7mm,below=9mm of ui] (r) {\textbf{Repositórios} \tiny fonte da verdade, cache, regras};
        \node[caixa,minimum width=1.8cm,below=3mm of r.south west,anchor=north west] (s1) {API\\ \tiny service};
        \node[caixa,minimum width=1.8cm,below=3mm of r.south] (s2) {Banco\\ \tiny service};
        \node[caixa,minimum width=1.8cm,below=3mm of r.south east,anchor=north east] (s3) {BLE\\ \tiny service};
        \node[pontilhado,fit=(r)(s1)(s2)(s3),label={[rotulo,anchor=south west]north west:\textbf{Camada de dados}}] (d) {};
        \draw[seta] (v) -- (vm);
        \draw[seta] (vm.south) -- (vm.south |- r.north);
        \draw[seta] (r.south -| s1) -- (s1); \draw[seta] (r) -- (s2); \draw[seta] (r.south -| s3) -- (s3);
      \end{tikzpicture}
    \end{column}
    \begin{column}{0.4\textwidth}
      \scriptsize
      \begin{tabularx}{\linewidth}{@{}l L@{}}
        \toprule
        \cabecalho{Peça} & \cabecalho{Responsabilidade} \\
        \midrule
        \textbf{View} & desenhar e repassar eventos; sem lógica \\
        \textbf{View model} & transformar dados em estado de tela e expor ações \\
        \textbf{Repositório} & fonte da verdade: cache, erros, retentativas \\
        \textbf{Serviço} & falar com uma fonte externa (HTTP, SQLite, BLE) \\
        \bottomrule
      \end{tabularx}
      \begin{dica}
        Uma view para cada view model. Repositórios e serviços são compartilhados entre telas.
      \end{dica}
    \end{column}
  \end{columns}
\end{frame}

\begin{frame}[fragile]{Organização das pastas}
  \begin{columns}[T]
    \begin{column}{0.46\textwidth}
\begin{codigo}[title=lib/,listing options={style=dm,language=bash,basicstyle=\ttfamily\tiny}]{bash}
lib/
|-- main.dart
|-- routing/router.dart
|-- domain/
|   `-- models/        bancada.dart, leitura.dart
|-- data/
|   |-- services/      api_service.dart, ble_service.dart
|   `-- repositories/  bancada_repository.dart
|-- ui/
|   |-- core/          tema e widgets comuns
|   |-- painel/
|   |   |-- view_models/painel_view_model.dart
|   |   `-- widgets/painel_screen.dart
|   `-- equipamentos/
`-- utils/             result.dart, command.dart
test/                  espelha lib/
\end{codigo}
    \end{column}
    \begin{column}{0.52\textwidth}
      \begin{conceito}[Organização mista]
        A camada de UI é organizada \textbf{por funcionalidade} (cada tela tem sua view e seu view model). A camada de dados é organizada \textbf{por tipo}, porque repositórios e serviços são usados por várias telas.
      \end{conceito}
      \begin{dica}[Referência oficial]
        O app de exemplo \textbf{Compass} (\cod{flutter/samples}) segue exatamente esta estrutura e é o estudo de caso da documentação.
      \end{dica}
    \end{column}
  \end{columns}
\end{frame}

% =============================================================================
\section{Camada de dados}

\begin{frame}[fragile]{Tratando falhas com \cod{Result}}
  \begin{columns}[T]
    \begin{column}{0.5\textwidth}
\begin{codigo}[title=utils/result.dart]{Dart}
sealed class Result<T> {
  const Result();
}

final class Ok<T> extends Result<T> {
  const Ok(this.valor);
  final T valor;
}

final class Falha<T> extends Result<T> {
  const Falha(this.erro);
  final Exception erro;
}
\end{codigo}
    \end{column}
    \begin{column}{0.48\textwidth}
\begin{codigo}[title=Usando]{Dart}
final r = await repo.listar();
switch (r) {
  case Ok(:final valor):
    bancadas = valor;
  case Falha(:final erro):
    mensagem = traduzir(erro);
}
\end{codigo}
      \begin{conceito}[Por que \cod{Result}?]
        O tipo de retorno deixa explícito que a operação pode falhar e o \cod{switch} exaustivo obriga quem chama a tratar o erro. Exceções esquecidas viram telas vermelhas.
      \end{conceito}
    \end{column}
  \end{columns}
\end{frame}

\begin{frame}[fragile]{Serviços e repositórios}
  \begin{columns}[T]
    \begin{column}{0.56\textwidth}
\begin{codigo}[title=data/repositories/bancada\_repository.dart]{Dart}
abstract class BancadaRepository {
  Future<Result<List<Bancada>>> listar();
}

class BancadaRepositoryRemoto implements BancadaRepository {
  BancadaRepositoryRemoto({required this.api});
  final ApiService api;
  List<Bancada>? _cache;

  @override
  Future<Result<List<Bancada>>> listar() async {
    if (_cache != null) return Ok(_cache!);
    try {
      final dtos = await api.getBancadas();   // serviço
      _cache = dtos.map((d) => d.paraModelo()).toList();
      return Ok(_cache!);
    } on Exception catch (e) { return Falha(e); }
  }
}
\end{codigo}
    \end{column}
    \begin{column}{0.42\textwidth}
      \begin{conceito}[Serviço]
        Envolve uma fonte externa e devolve dados brutos (DTOs). Não guarda estado nem decide nada.
      \end{conceito}
      \begin{conceito}[Repositório]
        Transforma DTOs em modelos de domínio, mantém cache, trata erros e decide de onde vêm os dados.
      \end{conceito}
      \begin{dica}
        A classe abstrata permite trocar a implementação (remota, local ou um \emph{fake} nos testes) sem mexer na UI.
      \end{dica}
    \end{column}
  \end{columns}
\end{frame}

% =============================================================================
\section{Camada de UI}

\begin{frame}[fragile]{O padrão Command}
  \begin{columns}[T]
    \begin{column}{0.54\textwidth}
\begin{codigo}[title=utils/command.dart (simplificado)]{Dart}
class Command0<T> extends ChangeNotifier {
  Command0(this._acao);
  final Future<Result<T>> Function() _acao;

  bool executando = false;
  Result<T>? resultado;
  bool get falhou => resultado is Falha;

  Future<void> execute() async {
    if (executando) return;          // evita duplo clique
    executando = true;
    resultado = null;
    notifyListeners();
    try {
      resultado = await _acao();
    } finally {
      executando = false;
      notifyListeners();
    }
  }
}
\end{codigo}
    \end{column}
    \begin{column}{0.44\textwidth}
      \begin{conceito}[Command]
        Envolve uma ação assíncrona do view model e expõe seu andamento: \textbf{executando}, \textbf{concluída} ou \textbf{falhou}. A view desenha cada estado sem \cod{bool}s espalhados.
      \end{conceito}
      \begin{dica}
        É o padrão usado no estudo de caso oficial do Flutter. Existe também a variante \cod{Command1<T, A>} para ações com argumento.
      \end{dica}
    \end{column}
  \end{columns}
\end{frame}

\begin{frame}[fragile]{O view model}
  \begin{columns}[T]
    \begin{column}{0.56\textwidth}
\begin{codigo}[title=ui/painel/view\_models/painel\_view\_model.dart]{Dart}
class PainelViewModel extends ChangeNotifier {
  PainelViewModel({required this.repo}) {
    carregar = Command0(_carregar)..execute();
  }
  final BancadaRepository repo;
  late final Command0<List<Bancada>> carregar;

  List<Bancada> _bancadas = [];  String _filtro = '';
  List<Bancada> get visiveis => _bancadas
      .where((b) => b.nome.contains(_filtro)).toList();
  void filtrar(String f) { _filtro = f; notifyListeners(); }

  Future<Result<List<Bancada>>> _carregar() async {
    final r = await repo.listar();
    if (r case Ok(:final valor)) _bancadas = valor;
    notifyListeners();
    return r;
  }
}
\end{codigo}
    \end{column}
    \begin{column}{0.42\textwidth}
      \small
      \begin{itemize}
        \item Recebe os repositórios no construtor (injeção).
        \item Guarda o estado da tela e o expõe só para leitura.
        \item Ações são métodos ou \emph{commands}.
        \item Não importa \cod{flutter/material.dart}: é Dart puro, fácil de testar.
      \end{itemize}
      \begin{atencao}
        O view model nunca guarda \cod{BuildContext} nem widgets.
      \end{atencao}
    \end{column}
  \end{columns}
\end{frame}

\begin{frame}[fragile]{A view}
  \begin{columns}[T]
    \begin{column}{0.58\textwidth}
\begin{codigo}[title=ui/painel/widgets/painel\_screen.dart]{Dart}
class PainelScreen extends StatelessWidget {
  const PainelScreen({super.key, required this.vm});
  final PainelViewModel vm;

  @override
  Widget build(BuildContext context) => ListenableBuilder(
    listenable: Listenable.merge([vm, vm.carregar]),
    builder: (context, _) {
      if (vm.carregar.executando) {
        return const Carregador();   // spinner centralizado
      }
      if (vm.carregar.falhou) {
        return ErroComRetentativa(
            onTentar: vm.carregar.execute);
      }
      return ListaBancadas(
          bancadas: vm.visiveis, onFiltro: vm.filtrar);
    },
  );
}
\end{codigo}
    \end{column}
    \begin{column}{0.4\textwidth}
      \begin{conceito}[Responsabilidade da view]
        Ler o estado do view model e desenhar; repassar eventos chamando métodos e \emph{commands}. Nenhuma regra de negócio.
      \end{conceito}
      \begin{dica}[Teste fácil]
        A tela recebe o view model no construtor: num teste, passe um view model com um repositório \emph{fake} (Módulo 14).
      \end{dica}
    \end{column}
  \end{columns}
\end{frame}

% =============================================================================
\section{Injeção de dependência}

\begin{frame}[fragile]{Injetando com Provider}
  \begin{columns}[T]
    \begin{column}{0.56\textwidth}
\begin{codigo}[title=main.dart]{Dart}
void main() {
  runApp(MultiProvider(
    providers: [
      Provider(create: (_) => ApiService()),
      Provider(create: (_) => BancoService()),
      Provider<BancadaRepository>(create: (ctx) =>
          BancadaRepositoryRemoto(api: ctx.read())),
      ChangeNotifierProvider(create: (_) => AuthRepository()),
    ],
    child: const LabMonitorApp(),
  ));
}

// na rota: cria o view model com as dependências
GoRoute(path: '/painel', builder: (ctx, _) =>
    PainelScreen(vm: PainelViewModel(repo: ctx.read())));
\end{codigo}
    \end{column}
    \begin{column}{0.42\textwidth}
      \begin{conceito}[Provider]
        Pacote recomendado pela documentação do Flutter para disponibilizar objetos na árvore de widgets. \cod{context.read<T>()} obtém o objeto; \cod{context.watch<T>()} também reconstrói quando ele notifica.
      \end{conceito}
      \begin{dica}
        Alternativas de DI: construtores (manual), \cod{get\_it} (\emph{service locator}) ou o próprio Riverpod.
      \end{dica}
    \end{column}
  \end{columns}
\end{frame}

% =============================================================================
\section{Gerência de estado: alternativas}

\begin{frame}[fragile]{Riverpod}
  \begin{columns}[T]
    \begin{column}{0.56\textwidth}
\begin{codigo}[title=Painel com Riverpod]{Dart}
final repoProvider = Provider<BancadaRepository>(
    (ref) => BancadaRepositoryRemoto(api: ref.watch(apiProvider)));

final bancadasProvider = FutureProvider<List<Bancada>>(
    (ref) => ref.watch(repoProvider).listarOuFalhar());

class PainelScreen extends ConsumerWidget {
  const PainelScreen({super.key});
  @override
  Widget build(BuildContext context, WidgetRef ref) {
    final estado = ref.watch(bancadasProvider);
    return estado.when(
      loading: () => const CircularProgressIndicator(),
      error: (e, _) => Text('Erro: $e'),
      data: (lista) => ListaBancadas(bancadas: lista),
    );
  }
}
// runApp(const ProviderScope(child: LabMonitorApp()));
\end{codigo}
    \end{column}
    \begin{column}{0.42\textwidth}
      \begin{conceito}[Riverpod]
        \emph{Providers} globais e tipados, independentes da árvore de widgets. Cache, carregamento e erro (\cod{AsyncValue}) já embutidos; dependências entre providers com \cod{ref.watch}.
      \end{conceito}
      \begin{dica}
        \cod{NotifierProvider} faz o papel do view model com estado mutável; geração de código com \cod{@riverpod} reduz repetição.
      \end{dica}
    \end{column}
  \end{columns}
\end{frame}

\begin{frame}[fragile]{Bloc / Cubit}
  \begin{columns}[T]
    \begin{column}{0.56\textwidth}
\begin{codigo}[title=Painel com Cubit]{Dart}
sealed class PainelEstado {}
class Carregando extends PainelEstado {}
class Pronto extends PainelEstado {
  Pronto(this.bancadas); final List<Bancada> bancadas; }
class ErroPainel extends PainelEstado {
  ErroPainel(this.msg); final String msg; }

class PainelCubit extends Cubit<PainelEstado> {
  PainelCubit(this._repo) : super(Carregando());
  final BancadaRepository _repo;

  Future<void> carregar() async {
    emit(Carregando());
    switch (await _repo.listar()) {
      case Ok(:final valor): emit(Pronto(valor));
      case Falha(:final erro): emit(ErroPainel('$erro'));
    }
  }
}
\end{codigo}
    \end{column}
    \begin{column}{0.42\textwidth}
\begin{codigo}[title=Na view]{Dart}
BlocProvider(
  create: (ctx) =>
      PainelCubit(ctx.read())..carregar(),
  child: BlocBuilder<PainelCubit,
      PainelEstado>(
    builder: (ctx, e) => switch (e) {
      Carregando() => const Carregador(),
      Pronto(:final bancadas) =>
          ListaBancadas(bancadas: bancadas),
      ErroPainel(:final msg) => Text(msg),
    }),
);
\end{codigo}
      \begin{conceito}[Bloc]
        Estados imutáveis emitidos em sequência; \cod{Bloc} completo recebe \textbf{eventos}, \cod{Cubit} recebe chamadas de método. Muito usado em equipes grandes.
      \end{conceito}
    \end{column}
  \end{columns}
\end{frame}

\begin{frame}{Qual escolher?}
  \footnotesize
  \begin{tabularx}{\textwidth}{@{}l L L L@{}}
    \toprule
    & \cabecalho{ChangeNotifier + Provider} & \cabecalho{Riverpod} & \cabecalho{Bloc / Cubit} \\
    \midrule
    Origem & SDK do Flutter + pacote recomendado & comunidade (autor do Provider) & comunidade (Very Good Ventures) \\
    Curva & baixa & média & média \\
    Boilerplate & baixo & baixo (com geração) & médio \\
    Pontos fortes & simples, oficial, MVVM direto & cache e assíncrono embutidos, testável & previsível, rastreável, padrão rígido \\
    Indicado para & começar e apps pequenos e médios & apps médios e grandes & equipes grandes, regras complexas \\
    \bottomrule
  \end{tabularx}
  \vspace{4pt}
  \begin{caixaambar}
    O que importa é a \dest{separação em camadas}: com ela, trocar a biblioteca de estado afeta só a camada de UI. Neste curso usamos \dest{ChangeNotifier + Provider}, como na documentação oficial.
  \end{caixaambar}
\end{frame}

% =============================================================================
\section{Fechamento}

\begin{frame}{Resumo}
  \begin{resumo}
    \begin{itemize}
      \item Duas camadas: UI (views + view models) e dados (repositórios + serviços).
      \item Views desenham; view models guardam estado e ações; repositórios são a fonte da verdade.
      \item \cod{Result} torna as falhas explícitas; \textbf{Command} expõe o andamento das ações.
      \item Dependências entram pelo construtor; Provider as disponibiliza na árvore.
      \item Riverpod e Bloc são alternativas maduras; a arquitetura em camadas vale para todas.
    \end{itemize}
  \end{resumo}
\end{frame}

\begin{frame}{Glossário (1/2)}
  \begin{glossario}
    Arquitetura & estruturas do sistema e decisões difíceis de mudar depois \\
    Camada & grupo de código com uma responsabilidade: UI ou dados (e domínio) \\
    MVVM & a \emph{view} observa o estado exposto pelo \emph{view model} \\
    MVI & estado único e imutável; cada intenção do usuário gera um novo estado \\
    \emph{View model} & transforma dados em estado de tela e expõe ações \\
    Repositório & fonte única da verdade de um tipo de dado \\
    Serviço & classe que fala com uma fonte externa: HTTP, SQLite, BLE \\
    \cod{Result} & tipo que representa sucesso (\cod{Ok}) ou falha (\cod{Falha}) sem exceções \\
    \emph{Command} & envolve uma ação e expõe seu andamento: executando, concluída, falhou \\
  \end{glossario}
\end{frame}

\begin{frame}{Glossário (2/2)}
  \begin{glossario}
    Regra da dependência & camadas internas não conhecem as externas \\
    Inversão de dependência & depender de abstrações, não de implementações \\
    Injeção de dependência & receber as dependências prontas em vez de criá-las \\
    Raiz de composição & único lugar onde os objetos concretos são criados e ligados \\
    Localizador de serviços & registro global ao qual a classe pede dependências (\cod{get\_it}) \\
    Provider e Riverpod & injeção e gerência de estado pela árvore ou por provedores \\
    Bloc e Cubit & estado emitido como fluxo, a partir de eventos ou de métodos \\
    ADR & registro curto de uma decisão de arquitetura e de seus motivos \\
  \end{glossario}
\end{frame}

\begin{frame}{Questões de revisão}
  \begin{questoes}
    \item Que problemas aparecem quando um widget faz HTTP, SQL e regras de negócio? Cite três.
    \item Compare MVP e MVVM: quem atualiza a tela em cada um?
    \item Qual a diferença de responsabilidade entre repositório e serviço?
    \item Enuncie a regra da dependência e dê um exemplo de violação.
    \item Como a interface \cod{FonteLeituras} permite trocar BLE por MQTT sem mexer no \emph{view model}?
    \item Qual a diferença entre injeção de dependência e localizador de serviços? Onde fica o \cod{get\_it}?
    \item Por que devolver \cod{Result} em vez de lançar exceções nos repositórios?
    \item Que critérios usar para escolher entre Provider, Riverpod e Bloc?
  \end{questoes}
\end{frame}

\begin{frame}{Exercícios}
  \begin{exercicio}[Exercício 1 --- Refatoração]
    Pegue uma tela com lógica dentro do widget e separe-a em view, view model, repositório e serviço.
  \end{exercicio}
  \begin{exercicio}[Exercício 2 --- Command]
    Implemente \cod{Command1<T, A>} e use-o na ação ``reservar bancada'', mostrando um indicador no botão enquanto executa e uma \emph{snackbar} em caso de falha.
  \end{exercicio}
  \begin{exercicio}[Exercício 3 --- Comparação]
    Implemente o mesmo painel com Riverpod e com Cubit. Compare quantidade de código, testabilidade e clareza.
  \end{exercicio}
  \begin{projeto}[LabMonitor --- Etapa 8]
    Reorganize o app nas pastas recomendadas: \cod{PainelViewModel}, \cod{EquipamentosViewModel}, \cod{BancadaRepository} com implementação \emph{fake} e injeção com Provider.
  \end{projeto}
\end{frame}

\begin{frame}{Para aprofundar}
  \small
  \begin{itemize}
    \item Guia de arquitetura: \url{https://docs.flutter.dev/app-architecture/guide}
    \item Estudo de caso (app Compass): \url{https://docs.flutter.dev/app-architecture/case-study}
    \item Recomendações de arquitetura: \url{https://docs.flutter.dev/app-architecture/recommendations}
    \item Provider: \url{https://pub.dev/packages/provider} \quad Riverpod: \url{https://riverpod.dev}
    \item Bloc: \url{https://bloclibrary.dev}
  \end{itemize}
\end{frame}

\begin{frame}{Respostas das questões de revisão}
  \begin{respostas}
    \item As regras não podem ser testadas sem montar a tela; uma mudança na API quebra várias telas; os dados ficam duplicados e inconsistentes; o arquivo fica enorme.
    \item No MVP, o \emph{presenter} chama métodos de uma interface da \emph{view}. No MVVM, o \emph{view model} expõe estado e a \emph{view} o observa, sem que ele a conheça.
    \item O serviço só conversa com uma fonte externa (HTTP, SQLite, BLE). O repositório combina serviços, decide cache e retentativas e é a fonte da verdade para a UI.
    \item As camadas internas não conhecem as externas. Violação: um repositório que importa \cod{package:flutter/material.dart} ou recebe um \cod{BuildContext}.
    \item O \emph{view model} depende só da interface; a raiz de composição injeta \cod{FonteBle} ou \cod{FonteMqtt}, e os testes usam uma \cod{FonteFake}.
    \item Na DI, a classe recebe as dependências de fora, explícitas no construtor. No localizador, ela mesma as pede a um registro global, como o \cod{get\_it}, o que as esconde e só falha em execução.
    \item O tipo de retorno deixa claro que a operação pode falhar e o \cod{switch} exaustivo obriga a tratar o erro; exceções esquecidas viram telas vermelhas.
    \item Tamanho do time e do app, familiaridade, necessidade de rastrear eventos (Bloc), verificação em compilação e testes (Riverpod) e simplicidade (Provider).
  \end{respostas}
\end{frame}

\encerramento{Módulo 09 (comum) --- Fundamentos de banco de dados: modelagem, SQL, normalização, transações e índices.}

\end{document}
